%!TEX root=../main.tex
\solutionFunctionAlt*
\begin{proof}
	Define
	\[
		\calX =
		\bigg\{ w \in \R^d :
		\wi = \alpha_\bi \ci, \alpha_\bi \geq 0, \forall \, \bi \in \equi, \,
		\wi = 0 \, \forall \, \bi \in \calB \setminus \equi,
		X w = \hat y.
		\bigg\}
	\]
	Let \( w \in \solfn(\lambda) \).
	If \( \wi \neq 0 \), then FO conditions imply
	\[
		\lambda \frac{\wi}{\norm{\wi}_2} = \ci \implies \wi = \alpha_\bi \ci,
	\]
	for \( \alpha > 0 \).
	If \( \wi = 0 \), then \( \wi = \alpha_\bi \ci \) holds trivially for
	\( \alpha_\bi = 0 \).
	Finally, since \( w \) is optimal, it must satisfy
	\[
		X \w = \hat y,
	\]
	which implies that \( w \in \calX \).

	Now, suppose that \( w' \in \calX \).
	Let's check FO conditions for \( w' \):
	if \( \wi \neq 0 \), then
	\begin{align*}
		\wi' = \alpha \ci
		 & \implies \frac{\wi'}{\norm{\wi'}_2} = \frac{\ci}{\lambda} \\
		 & \implies \lambda \frac{\wi'}{\norm{\wi'}_2} = \ci,
	\end{align*}
	where we have used the fact that \( \norm{\ci}_2 = \lambda \)
	for all \( \bi \in \equi \).
	If \( \wi' = 0 \), then
	\[
		X w' = \hat y \implies \norm{\Xbi^\top (y - Xw')}_2 \leq \lambda,
	\]
	which satisfies FO conditions.
	We conclude \( w' \in \solfn(\lambda) \).
\end{proof}


\begin{restatable}{lemma}{boundedSolutions}\label{lemma:bounded-solutions}
	The correlation vector \( v(\lambda) = X^\top (y -  X\wmin(\lambda)) \) is bounded
	for all \( \lambda \geq 0 \).
	Moreover, every solution to the group lasso problem
	is bounded by an absolute constant independent of \( \lambda \).
	Specifically, every \( w \in \solfn(\lambda) \) satisfies
	\[
		\sum_{\bi \in \calB} \norm{\wi}_2 \leq \sum_{\bi \in \calB} \norm{\bwi}_2.
	\]
	where \( \bar w \) is the min-norm least-squares solution.
\end{restatable}
\begin{proof}
	As noted, for \( \lambda = 0 \), the min-norm solution reduces to the min-norm least-squares
	solution and \( v(0) = 0 \) so that the claim trivial holds.
	For \( \lambda > 0 \), we show boundedness directly.
	Optimality of \( \wmin \) implies
	\[
		\half \norm{X \wmin - y}_2^2 + \lambda \sum_{\bi \in \calB} \norm{\wmin_\bi}_2
		\leq \half \norm{y}_2^2,
	\]
	from which we deduce the residual vector \( r(\lambda) \) is bounded
	for any choice of \( \lambda \).
	Thus, the block correlation vectors satisfy,
	\begin{align*}
		\norm{\vi(\lambda)}^2_2      & = \norm{\Xbi^\top r(\lambda)}_2^2                                   \\
		                             & \leq \norm{\Xbi}_2^2 \norm{y}_2^2                                   \\
		\implies \norm{v(\lambda)}_2 & \leq \sbr{\sum_{\bi \in \calB} \norm{\Xbi}^2_2 \norm{y}^2_2}^{1/2}.
	\end{align*}
	The right-hand quantity is independent of \( \lambda \) and
	so \( v(\lambda) \) is bounded for all \( \lambda \).

	Let \( h(w) = \sum_{\bi \in \calB} \norm{\wi}_2 \),
	and define \( W_{g} \) to be the set of least squares solutions with
	minimum group norm.
	That is,
	\[
		W_{g} = \argmin \cbr{h(w) : X^\top X w = X^\top y }.
	\]
	Let \( w_g \in W_g \),
	and suppose that \( h(w_{g}) < h(w(\lambda)) \)
	for some \( \lambda > 0 \), \( w \in \solfn(\lambda) \).
	Since
	\[
		\min_{w} \half \norm{X w - y}_2^2
		\leq \min_{w} \half \norm{X w - y}_2^2
		+ \lambda h(w),
	\]
	we deduce
	\[
		\half \norm{X w_g - y}_2^2 + \lambda h(w_g) <
		\half \norm{X \w(\lambda) - y}_2^2
		+ \lambda h(\w(\lambda)),
	\]
	which is a contradiction.
	So \( h(\w(\lambda)) \leq h(w_g) \) for all \( \lambda > 0 \).
	Observing \( h(w_g) \leq h(\bar w) \) by definition gives the result.
	Since \( h(\bar w) \) is independent of \( \lambda \),
	we conclude that \( \wmin(\lambda) \) is bounded.
\end{proof}


\subsection{Properties of the Solution Map}\label{app:solution-map}

\valueContinuity*
\begin{proof}
	Define the joint objective function
	\[
		f(w, \lambda) = \half \norm{X w - y}_2^2 + \lambda \sum_{b_i \in \calB} \norm{\wi}_2.
	\]
	Clearly \( f(w, \lambda) \) is jointly continuous in \( w \) and \( \lambda \).
	By \cref{lemma:bounded-solutions}, minimization of
	\( f(w, \lambda) \) over \( \R^d \) is equivalent to the
	constrained minimization problem,
	\[
		p^*(\lambda) = \min_{w} f(w, \lambda) \quad \text{s.t.} \; \sum_{\bi \in \calB} \norm{\wi}_2 \leq C,
	\]
	where \( C \) is a finite absolute constant.
	Note that this expression is also valid when \( \lambda = 0 \) as the
	min-norm solution to the unregularized least squares problem
	obeys the constraint.
	Thus, \( \solfn \) is a continuous optimization problem over a compact
	set and the classical result of \citet{berge1997topological}
	(see also \citet{hogan1973point}[Theorem 7]) implies
	\( p^* \) is continuous.
\end{proof}



\mapContinuity*
\begin{proof}
	The result follows from joint continuity of the objective
	\[
		f(w, \lambda) = \half \norm{X w - y}_2^2 + \lambda \sum_{b_i \in \calB} \norm{\wi}_2,
	\]
	and straightforward application of \citet{robinson1974sufficient}[Theorem 1].

	If \( X \) is full column rank, then the group lasso solution is
	unique for all \( \lambda \geq 0 \).
	The solution map is a singleton on \( \R^+ \) and closedness
	and openness are equivalent properties for singleton maps.
	Since we have already shown it is closed, the solution map
	must also be open.
	An identical argument shows that the solution map is open
	at every \( \lambda > 0 \) when GGP is satisfied.

	Now we show the reverse implication by proving the contrapositive.
	Assume \( X \) is not full column-rank.
	The solution map at \( \lambda = 0 \) is the solution set to the
	(unregularized) least squares problem,
	\[
		\min_w \half \norm{X w - y}_2^2,
	\]
	which is known to be \( \solfn(0) = \cbr{\wmin(0) + z : z \in \Null(X)} \).
	While \( \solfn(0) \) is unbounded, it holds that
	\( \solfn(\lambda) \subset C \) for some bounded \( C \) for
	every \( \lambda > 0 \) (\cref{lemma:bounded-solutions}).
	As a result, there exist uncountably many solutions in \( \solfn(0) \)
	which are not limit points of solutions in \( \solfn(\lambda_k) \)
	as \( \lambda_k \rightarrow 0 \).
	In other words, \( \solfn(0) \) is not open at \( 0 \).

	%Suppose \( X \) is not full column rank.
	%The solution map \( \solfn \) is open if and only if it is lower semi-continuous
	%\citep{hogan1973point}.
	%Recall that a point-to-set map \( T \) is lower semi-continuous at \( x \) if
	%for every open set \( S \), \( T(x) \cap S \neq \emptyset \), there
	%exists a neighborhood \( N(x) \) of \( x \) such that for all \( x' \in N(x) \),
	%\( T(x') \cap S \neq \emptyset \).

	%Suppose that \( \solfn \) is open from the right at \( \lambda = 0 \).
	%Standard results imply \( \text{aff}(\solfn(0)) = \Null(X) + w \),
	%where \( w \in \solfn(0) \).
	%We shall assume that \( w \) is dense.
	%By openness, for every \( \lambda_k \into 0 \), there exists
	%\( k' \),
	%and \( w(\lambda_k) \in \solfn(\lambda_k) \)
	%for \( k \geq k' \)
	%such that \( w(\lambda_k) \rightarrow w \).
	%Then \( w(\lambda_k) \) is also dense for sufficiently
	%large \( k \) and \( \calE_{\lambda_k} = \calB \).
	%\cref{cor:uniqueness-cond} now implies
	%\( \text{aff}(\solfn(\lambda')) \subset \Null(X) \) for all
	%\( \lambda' \) in a sufficiently small neighbourhood of
	%\( 0 \).

	%Since \( \solfn \) is open from the right at \( 0 \),
	%it is lower semi-continuous from the right.
	%Let
	%\[
	%    S = \text{relint}(\solfn(0))
	%    + B(0, \epsilon) \cap \Row(X),
	%\]
	%which is clearly open and satisfies
	%\( S \cap \solfn(0) = \text{relint}(\solfn(\lambda)) \),
	%which is always non-empty.
	%By lower semi-continuity, there exists a neighborhood
	%\( N(0) \) of \( 0 \) such that for every \( \lambda' \in N(0) \),
	%\( \solfn(\lambda') \cap S \neq 0 \).
	%That is, there exists \( w_1 \in \text{relint}(\solfn(0)), z_1 \in B(0, \epsilon) \cap \Row(X) \),
	%and \( W^{(2)} \in \solfn(\lambda') \)
	%such that
	%\[
	%    w_1 + z_1 = W^{(2)}.
	%\]
	%But \( W^{(2)} \) must be orthogonal to \( z_1 \)
	%since \( \text{aff}(W^{(2)}) \subset \Null(X) \).
	%Since \( w_1 \) is also orthogonal to \( z_1 \),
	%we conclude \( z_1 = 0 \) and \( w_1 = W^{(2)} \).

	%However, this condition cannot happen because
	%\[
	%    \norm{\vi(0)}_2 = 0 \neq \lambda' = \norm{\vi(\lambda')}_2.
	%\]
	%So, the solution map is not open from the right at \( \lambda = 0 \).
\end{proof}


